%% kapa: exam=1 round=32 year=2021 period=2 form=A subject=appraisal_law subject_ko=감정평가관계법규 seq=24 number=24
%% source: Q-Net 감정평가사 기출문제 (cst00309) artlSeq=5208397 file=extracted/1008/5208397/A형/2021년도 제32회 감정평가사 1차 2교시 A형.pdf
%% conversion: kordoc 4.24.51 -> tools/gen1.py (auto); review=proofread-sonnet; flags=none
\begin{question}{24}
건축법령상 용어에 관한 설명으로 옳지 않은 것은?
\begin{choices}
  \item “지하층”이란 건축물의 바닥이 지표면 아래에 있는 층으로서 바닥에서 지표면까지 평균높이가 해당 층 높이의 3분의 1 이상인 것을 말한다.
  \item “거실”이란 건축물 안에서 거주, 집무, 작업, 집회, 오락, 그 밖에 이와 유사한 목적을 위하여 사용되는 방을 말한다.
  \item “고층건축물”이란 층수가 30층 이상이거나 높이가 120미터 이상인 건축물을 말한다.
  \item “초고층 건축물”이란 층수가 50층 이상이거나 높이가 200미터 이상인 건축물을 말한다.
  \item “이전”이란 건축물의 주요구조부를 해체하지 아니하고 같은 대지의 다른 위치로 옮기는 것을 말한다.
\end{choices}
\end{question}
\answer{1}
